%
% reconciliation_hub_reject.tex
%
% A Z specification of the CONVERGENCE half of the lux commit-on-idle
% reconciliation: what happens when the value a field commits is one the Hub
% will REJECT rather than install.  It is a focused sibling of
% input_text_reconciliation.tex, which models the honour / defer / commit-echo
% discipline over a carrier [VALUE] that is IMPLICITLY the set of Hub-accepted
% values.  That implicit assumption -- every committed value is one the Hub
% installs -- is the one this document makes explicit and then defends.
%
% The defect class (honour-forever divergence) surfaced twice on input_number:
%   Round 1: a typed out-of-range value (150 into [0,100]) was committed; the
%     Hub's range check rejected it.
%   Round 2: a non-finite value (1e400 -> +inf) on a field unbounded on that
%     side clamped through unchanged and was committed; the Hub's finiteness
%     check rejected it.
% In both, the renderer committed a value OUTSIDE the Hub-accepted domain.  The
% Hub's apply_patch raises ValueError; EventHandlerHost.fire swallows it and does
% not re-push; so the Hub value never advances to the committed value.  The
% arbiter forgets its optimistic echo only when the Hub value moves OFF the
% commit-time value -- which now never happens -- so the Display honours the
% rejected value FOREVER while the Hub sits on the old one.  A silent, permanent
% Hub/Display divergence: no test round finds it, because no test drives the
% renderer past the Hub's boundary.
%
% This model distinguishes the Hub-accepted subset ``accepted'' from an
% arbitrary committed value, adds an explicit Hub-REJECT transition, and states
% the convergence invariant the defect violates.  With the reject transition
% enabled and Commit UNGUARDED, ProB reaches the honour-forever state; with the
% fix -- Commit guarded so it only ever commits an accepted value -- the reject
% transition is dead and convergence holds.
%
% Source of record:
%   src/punt_lux/protocol/elements/input_number.py
%       -- InputNumberElement.apply_patch: re-checks range AND finiteness at the
%          boundary and raises ValueError on any invalid result (the REJECT);
%          .clamped(): projects into [min,max] but does NOT sanitise non-finite,
%          so +inf on an unbounded side passes through -- the round-2 hole.
%       -- NumericInputChecks: the range/finiteness/format predicate defining the
%          Hub-accepted domain (this model's ``accepted'').
%   src/punt_lux/display/renderers/input_number_renderer.py
%       -- InputNumberRenderer.render / ._commit: projects the raw entry through
%          elem.clamped before commit.  The fix extends that projection to a total
%          sanitiser into the accepted domain, so ``committed in accepted'' holds
%          by construction -- this model's Commit guard.
%   src/punt_lux/display/renderers/imgui/continuous_edit_selection.py
%       -- ContinuousEditArbiter.resolve: returns the committed value while the
%          Hub value still equals the commit-time Hub value; forgets it once the
%          Hub value moves off -- the branch that never fires under a rejected
%          commit.
%
% Type-check:  fuzz -t docs/reconciliation_hub_reject.tex
% Model-check (see Section "Verification" for the exact goals):
%   probcli docs/reconciliation_hub_reject.tex -model_check -p DEFAULT_SETSIZE 2
%   probcli docs/reconciliation_hub_reject.tex -model_check -nodead \
%       -goal "pending=ztrue & committed /: accepted" -p DEFAULT_SETSIZE 2
%   probcli docs/reconciliation_hub_reject.tex -model_check -nodead \
%       -goal "diverged=ztrue" -p DEFAULT_SETSIZE 2
%   probcli docs/reconciliation_hub_reject.tex -model_check -nodead \
%       -goal "editing=ztrue & edited=zfalse" -p DEFAULT_SETSIZE 2
%
% Formatting follows the fuzz house rules: \quad~ for continuation indent,
% schema lines under 80 columns, predicates broken at \land/\lor/\implies with
% the operator starting the continuation line.
%

\documentclass[a4paper,10pt,fleqn]{article}
\usepackage[margin=1in]{geometry}
\usepackage{fuzz}
\usepackage[colorlinks=true,linkcolor=blue,citecolor=blue,urlcolor=blue]{hyperref}

\begin{document}

\title{Hub Rejection and Convergence in the lux Commit-on-Idle Discipline:
  a Z Specification}
\author{Formal model of the honour-forever divergence and its exclusion}
\date{July 2026}
\maketitle

% ============================================================
\section{Introduction}
% ============================================================

The commit-on-idle reconciliation of \texttt{input\_text\_reconciliation.tex}
governs a field whose value is Hub-authoritative and replicated to a Display.
Its carrier \texttt{[VALUE]} stands for whatever a field holds --- text, a float,
an RGBA tuple --- and its \texttt{HubEcho} operation installs the committed value
at the Hub unconditionally.  That unconditional install encodes an assumption the
present document makes explicit: \emph{every value a field commits is one the Hub
will accept}.  The carrier is, in effect, the set of Hub-accepted values.

The assumption is false for a numeric field.  An \texttt{input\_float} widget does
not clamp, and the renderer's own projection into range, \texttt{elem.clamped},
leaves a value non-finite when the field is unbounded on the offending side: a
typed \texttt{1e400} becomes \texttt{+inf} and passes through.  The field commits
that value; the Hub's \texttt{apply\_patch} re-checks range \emph{and} finiteness
and raises \texttt{ValueError}; the firing host logs the error and does not
re-push.  The Hub value therefore never advances to the committed value.

The arbiter's optimistic echo is what turns a rejected commit into a permanent
fault.  On commit the arbiter records the committed value and the Hub value
observed at commit time, and while the field is idle it returns the committed
value until the Hub value moves \emph{off} that commit-time value --- the moment
it treats as the echo landing.  When the Hub rejects the commit, the Hub value
never moves; the arbiter's forget branch never fires; and the Display honours the
rejected value for the rest of the field's life while the Hub holds the old one.
This is the \emph{honour-forever divergence}.  It surfaced twice on
\texttt{input\_number} --- once for an out-of-range value, once for a non-finite
one --- and that recurrence is why the boundary is modelled here rather than
patched a third time.

The correction is a single discipline: \emph{a field commits only a value the Hub
will accept}.  The renderer's pre-commit projection is widened from a range clamp
to a total sanitiser onto the accepted domain, so a non-finite or out-of-range
entry is repaired before it is committed, and the value that travels to the Hub is
always one \texttt{apply\_patch} installs.  The Hub then always echoes, the forget
branch always eventually fires, and the display converges.

This document distinguishes the Hub-accepted subset from an arbitrary committed
value, adds an explicit Hub-reject transition, and states the convergence
property the defect violates.  Over a small bounded carrier ProB enumerates every
interleaving of edit, commit, echo, and reject, and either confirms convergence or
produces the exact honour-forever trace.

% ============================================================
\section{Basic Types}
% ============================================================

\begin{zed}
  [VALUE]
\end{zed}

\texttt{VALUE} is the set of values the field may hold and the widget may commit.
Crucially it is \emph{not} restricted to Hub-accepted values: the widget can
produce a value the Hub refuses, and the whole point of this model is that such a
value can be committed.  Two values suffice to exhibit the divergence --- an
accepted starting value and one value the Hub rejects; three widen the
interleaving without adding a reachable fault.  The carrier is bounded by ProB's
\texttt{DEFAULT\_SETSIZE}.

% ============================================================
\section{Free Types and Constants}
% ============================================================

\begin{zed}
  ZBOOL ::= ztrue | zfalse
\end{zed}

A two-valued flag written as a free type rather than the toolkit \texttt{bool}, so
ProB animates it as a small enumerated set.

\begin{axdef}
  v0 : VALUE \\
  accepted : \power VALUE
\where
  v0 \in accepted
\end{axdef}

\texttt{accepted} is the Hub-accepted domain: the set of values
\texttt{apply\_patch} installs rather than rejecting.  For a numeric field it is
the finite values within \texttt{[min, max]}; the model leaves it an arbitrary
subset so ProB checks every shape of boundary, including the one where only the
start value is accepted.  \texttt{v0} is that start value --- both the initial
display and the initial Hub value --- and it must be accepted, since the Hub holds
it.  A value in \texttt{VALUE} but outside \texttt{accepted} is one the widget can
commit and the Hub will reject; with \texttt{DEFAULT\_SETSIZE 2} and
\texttt{accepted} taken as \texttt{\{v0\}}, the other value is exactly such a one.

% ============================================================
\section{State}
% ============================================================

The state is flat: one schema, eight components.  Its predicate carries only
structural coherence --- that a field can defer or be mid-edit only while focused.
That holds in every design, correct or buggy, so it constrains rather than tests.
The convergence property (Section~\ref{sec:inv}) is deliberately \emph{absent}
here, so a divergent state stays reachable and ProB can find it.

\begin{schema}{Field}
  disp : VALUE \\
  hub : VALUE \\
  focused : ZBOOL \\
  editing : ZBOOL \\
  edited : ZBOOL \\
  pending : ZBOOL \\
  committed : VALUE \\
  diverged : ZBOOL
\where
  editing = ztrue \implies focused = ztrue \\
  edited = ztrue \implies focused = ztrue
\end{schema}

\texttt{disp} is the value the Display renders --- the arbiter's \texttt{resolve}
output.  \texttt{hub} is the Hub-authoritative value.  \texttt{focused},
\texttt{editing} (the defer flag) and \texttt{edited} (the real-edit witness)
carry the honour/defer discipline unchanged from the parent model, kept so the
convergence question is asked over a faithful edit cycle rather than an abstract
one.

\texttt{pending} marks the echo-latency window: set from a commit, and cleared
only when the Hub adopts the committed value.  While it holds, \texttt{resolve}
returns \texttt{committed}, so the display shows the committed value.
\texttt{committed} is the value most recently committed --- possibly one outside
\texttt{accepted}.  \texttt{diverged} is pure bookkeeping for the invariant: it
latches \texttt{ztrue} the moment the Hub rejects an outstanding commit, the event
after which the display honours a value the Hub will never hold.

% ============================================================
\section{Initialization}
% ============================================================

\begin{schema}{Init}
  Field'
\where
  disp' = v0 \\
  hub' = v0 \\
  focused' = zfalse \\
  editing' = zfalse \\
  edited' = zfalse \\
  pending' = zfalse \\
  committed' = v0 \\
  diverged' = zfalse
\end{schema}

A single initial state: the field unfocused and idle, display and Hub both on the
accepted value \texttt{v0}, no commit outstanding, nothing diverged.

% ============================================================
\section{The Edit Cycle}
% ============================================================

Focusing begins a fresh session and, following the corrected parent design, does
\emph{not} begin deferring: a focused-but-unedited field still honours the Hub.

\begin{schema}{Focus}
  \Delta Field
\where
  focused = zfalse \\
  focused' = ztrue \\
  editing' = zfalse \\
  edited' = zfalse \\
  disp' = disp \\
  hub' = hub \\
  pending' = pending \\
  committed' = committed \\
  diverged' = diverged
\end{schema}

A keystroke changes the text and begins deferring.  The typed value \texttt{e?}
ranges over the whole of \texttt{VALUE} --- including values outside
\texttt{accepted}: this is where the user types \texttt{1e400}, the value the Hub
will later reject.  Modelling \texttt{e?} as unrestricted is what lets the widget
present the Hub with a value it refuses.

\begin{schema}{Keystroke}
  \Delta Field \\
  e? : VALUE
\where
  focused = ztrue \\
  editing' = ztrue \\
  edited' = ztrue \\
  disp' = e? \\
  focused' = focused \\
  hub' = hub \\
  pending' = pending \\
  committed' = committed \\
  diverged' = diverged
\end{schema}

A commit is the deactivation of an edited field.  It records the committed value
and opens the echo-latency window; the Hub is not updated here.  The guard
\texttt{disp $\in$ accepted} is the \emph{fix}: the field commits only a value the
Hub will accept.  It models the renderer's total pre-commit sanitiser --- a
non-finite or out-of-range entry is repaired into \texttt{accepted} before this
point, so the committed value is always installable.  Deleting the guard is the
defect (Section~\ref{sec:fidelity}): the widget then commits whatever it displays,
rejected values included.

\begin{schema}{Commit}
  \Delta Field
\where
  focused = ztrue \\
  editing = ztrue \\
  disp \in accepted \\
  committed' = disp \\
  pending' = ztrue \\
  focused' = zfalse \\
  editing' = zfalse \\
  edited' = zfalse \\
  disp' = disp \\
  hub' = hub \\
  diverged' = diverged
\end{schema}

A field may also lose focus without an edit.  This commits nothing.

\begin{schema}{BlurNoEdit}
  \Delta Field
\where
  focused = ztrue \\
  editing = zfalse \\
  focused' = zfalse \\
  edited' = zfalse \\
  editing' = zfalse \\
  disp' = disp \\
  hub' = hub \\
  pending' = pending \\
  committed' = committed \\
  diverged' = diverged
\end{schema}

% ============================================================
\section{The Hub Response}
% ============================================================

A committed value reaches the Hub, which either installs it or rejects it.  The
two outcomes are the two operations here, and they partition on whether the
committed value is accepted.

\texttt{HubEchoAccept} is the delayed echo of an \emph{accepted} commit: the Hub
adopts the committed value and the echo-latency window closes.  From here an
honour frame renders \texttt{hub}, which now equals \texttt{committed}: the
optimistic display and the real Hub value have converged.

\begin{schema}{HubEchoAccept}
  \Delta Field
\where
  pending = ztrue \\
  committed \in accepted \\
  hub' = committed \\
  pending' = zfalse \\
  disp' = disp \\
  focused' = focused \\
  editing' = editing \\
  edited' = edited \\
  committed' = committed \\
  diverged' = diverged
\end{schema}

\texttt{HubReject} is the delayed \emph{refusal} of a commit the Hub will not
install: the committed value is outside \texttt{accepted}, so \texttt{apply\_patch}
raises and the firing host swallows the error.  The Hub value does not move, and
the echo-latency window does not close --- \texttt{pending} stays set, so
\texttt{resolve} keeps returning the rejected value.  The operation latches
\texttt{diverged}: from this state the display honours a value the Hub will never
hold, and no operation can carry the Hub to it, because only
\texttt{HubEchoAccept} advances \texttt{hub} and it is disabled while the committed
value is unaccepted.

\begin{schema}{HubReject}
  \Delta Field
\where
  pending = ztrue \\
  committed \notin accepted \\
  diverged' = ztrue \\
  hub' = hub \\
  pending' = pending \\
  disp' = disp \\
  focused' = focused \\
  editing' = editing \\
  edited' = edited \\
  committed' = committed
\end{schema}

% ============================================================
\section{The Render Frames}
\label{sec:frames}
% ============================================================

A frame in which the field is not deferring honours: it drives the display to the
value the field is authoritative for --- the committed value while a commit is in
flight, the raw Hub value otherwise.  This is the \texttt{resolve} rule: it is
what makes the display show the committed value throughout the window, and hence
what makes a rejected commit visible for as long as the window stays open.

\begin{schema}{IdleHonour}
  \Delta Field
\where
  focused = zfalse \\
  editing = zfalse \\
  disp' = \IF pending = ztrue \THEN committed \ELSE hub \\
  focused' = focused \\
  editing' = editing \\
  edited' = edited \\
  hub' = hub \\
  pending' = pending \\
  committed' = committed \\
  diverged' = diverged
\end{schema}

\begin{schema}{FocusedHonour}
  \Delta Field
\where
  focused = ztrue \\
  editing = zfalse \\
  disp' = \IF pending = ztrue \THEN committed \ELSE hub \\
  focused' = focused \\
  editing' = editing \\
  edited' = edited \\
  hub' = hub \\
  pending' = pending \\
  committed' = committed \\
  diverged' = diverged
\end{schema}

A deferring field --- \texttt{editing} set --- honours nothing: its display is the
buffer, moved only by a further keystroke.  That case is the absence of an honour
frame while \texttt{editing} holds.  It is what keeps a step always available: an
unfocused field can \texttt{IdleHonour}; a focused non-deferring field can
\texttt{FocusedHonour}; a deferring field is focused, so it can \texttt{Keystroke}.
A step is therefore enabled in every reachable state, and the machine cannot
deadlock (Section~\ref{sec:inv}, invariant~3).

% ============================================================
\section{Invariants}
\label{sec:inv}
% ============================================================

Three properties must hold across all reachable states.  None is placed in the
\texttt{Field} schema: a predicate placed there would guard every successor,
silently disabling any operation that would break it and masking the very defect
the model exists to catch.

\paragraph{Invariant 1 --- convergence (no honour-forever).}
Whenever a commit is outstanding, the outstanding value is one the Hub will
accept.  Then the echo can always land: \texttt{HubEchoAccept} is enabled, the
window closes, and the display converges to the Hub.  The negation --- a commit
outstanding whose value the Hub will refuse --- is the honour-forever state, in
which \texttt{resolve} returns the committed value on every frame while no
operation can carry the Hub to it.  The invariant is that this state is never
reached:
\[
  pending = ztrue \implies committed \in accepted.
\]

\paragraph{Invariant 2 --- rejection never occurs.}
Equivalently, and as a named event: the Hub never rejects an outstanding commit,
so the divergence latch never fires:
\[
  diverged = zfalse.
\]
The two forms coincide.  \texttt{HubReject} is enabled exactly when invariant~1's
negation holds, so if invariant~1 holds on every reachable state then
\texttt{HubReject} never fires and \texttt{diverged} stays clear; and if
invariant~1's negation is reachable then \texttt{HubReject} fires from that state
and latches \texttt{diverged}.  Invariant~1 is the property; invariant~2 is the
witness that makes the offending trace end on a single named transition.

\paragraph{Invariant 3 --- deadlock freedom.}
No reachable state leaves the machine unable to progress.  The honour frames cover
the not-deferring case; a deferring field is focused and can \texttt{Keystroke};
so a step is always enabled.

% ============================================================
\section{Verification}
\label{sec:verify}
% ============================================================

Type-checking is a \texttt{fuzz} pass:
\begin{verbatim}
  fuzz -t docs/reconciliation_hub_reject.tex
\end{verbatim}

Invariants 1 and 2 are state properties, checked by asking ProB whether their
negation is reachable.  A goal not found means the invariant holds on every
reachable state:
\begin{verbatim}
  # Invariant 1 -- convergence (must report: goal NOT found)
  probcli docs/reconciliation_hub_reject.tex -model_check -nodead \
    -goal "pending=ztrue & committed /: accepted" -p DEFAULT_SETSIZE 2

  # Invariant 2 -- rejection never occurs (must report: goal NOT found)
  probcli docs/reconciliation_hub_reject.tex -model_check -nodead \
    -goal "diverged=ztrue" -p DEFAULT_SETSIZE 2

  # Inherited coherence -- defer only after a real edit (must report: NOT found)
  probcli docs/reconciliation_hub_reject.tex -model_check -nodead \
    -goal "editing=ztrue & edited=zfalse" -p DEFAULT_SETSIZE 2
\end{verbatim}

Invariant 3 is the deadlock check over the full reachable state space:
\begin{verbatim}
  # Invariant 3 -- deadlock freedom (must report: no deadlock, no invariant viol.)
  probcli docs/reconciliation_hub_reject.tex -model_check -p DEFAULT_SETSIZE 2
\end{verbatim}

All return the same verdict at \texttt{DEFAULT\_SETSIZE 3}: the extra value widens
the interleaving of a second edit against an in-flight commit without adding a
reachable violation.  Because \texttt{accepted} is left underspecified, ProB
discharges each goal over \emph{every} valuation of the accepted domain that
contains \texttt{v0} --- so convergence is proved not for one boundary but for all.

% ============================================================
\section{Fidelity: reproducing the honour-forever divergence}
\label{sec:fidelity}
% ============================================================

A model that cannot reproduce the known defect proves nothing.  The defect is the
corrected specification with the commit guard removed: strike
\texttt{disp $\in$ accepted} from \texttt{Commit}, so the field commits whatever it
displays, rejected values included.  Nothing else changes.  Under
\texttt{DEFAULT\_SETSIZE 2}, take \texttt{accepted = \{v0\}}, and write $v_0 = v0$
for the accepted start value and $v_r$ for the other, rejected, value.

With the guard gone, both convergence goals turn from \emph{not found} to
\emph{found}.  The shortest trace to \texttt{diverged = ztrue}:
\begin{enumerate}
  \item \texttt{Init}: \texttt{disp = hub = committed = $v_0$}, unfocused, idle,
    \texttt{pending = zfalse}, \texttt{diverged = zfalse}.
  \item \texttt{Focus}: \texttt{focused = ztrue}, still not editing.
  \item \texttt{Keystroke($v_r$)}: the user types the value the Hub will reject.
    \texttt{disp = $v_r$}, \texttt{editing = ztrue}.
  \item \texttt{Commit} (unguarded): \texttt{committed = $v_r$},
    \texttt{pending = ztrue}, \texttt{hub = $v_0$}, \texttt{focused = zfalse},
    \texttt{editing = zfalse}.  A commit is now outstanding for a value outside
    \texttt{accepted} --- invariant~1's negation already holds:
    \texttt{pending = ztrue} and \texttt{committed = $v_r$ $\notin$ accepted}.
  \item \texttt{IdleHonour}: \texttt{pending} holds, so \texttt{disp' = committed
    = $v_r$}.  The Display shows the rejected value.
  \item \texttt{HubReject}: \texttt{committed = $v_r$ $\notin$ accepted}, so the
    Hub refuses.  \texttt{hub} stays $v_0$, \texttt{pending} stays \texttt{ztrue},
    and \texttt{diverged} latches \texttt{ztrue}.
\end{enumerate}
From the state after step~6 the divergence is permanent.  \texttt{HubEchoAccept}
is disabled (\texttt{committed $\notin$ accepted}); \texttt{HubReject} only
re-latches; every \texttt{IdleHonour} rewrites \texttt{disp = committed = $v_r$}.
The Display honours $v_r$ forever, the Hub holds $v_0$ forever, and the arbiter's
forget branch --- which fires only when \texttt{hub} moves off the commit-time
value --- never fires because \texttt{hub} never moves.  This is exactly the
field-observed fault: a value typed past the Hub's boundary, committed, echoed
nowhere, and displayed indefinitely.

Under the corrected design the guard \texttt{disp $\in$ accepted} makes step~4
impossible for $v_r$: \texttt{Commit} is disabled while the display holds a
rejected value, so the user must repair the entry (a further \texttt{Keystroke}
into \texttt{accepted}) or blur.  \texttt{committed $\in$ accepted} therefore holds
throughout every window; \texttt{HubReject} is dead; \texttt{HubEchoAccept} always
eventually fires; and the reachability search confirms
\texttt{pending = ztrue \& committed /: accepted} and \texttt{diverged = ztrue} are
both \emph{not found}, for every valuation of \texttt{accepted}.

% ============================================================
\section{From the Model to the Code}
% ============================================================

The model's commit guard is one change to the numeric renderer.  Today
\texttt{InputNumberRenderer} projects a raw entry through \texttt{elem.clamped}
before committing; \texttt{clamped} bounds the value into \texttt{[min, max]} but
does nothing to a non-finite value on an unbounded side, so \texttt{+inf} survives
to the commit and the Hub rejects it.  The corrected projection is a \emph{total
sanitiser} onto the accepted domain: it maps every widget value --- non-finite
included --- to a finite in-range value before the value is observed, committed, or
fired.  A non-finite entry collapses to a bound (or to the field's current value
when the field is unbounded), an out-of-range entry clamps, and the value that
reaches the Hub is one \texttt{apply\_patch} installs without raising.

Concretely, the pre-commit projection must satisfy, for every widget value
\texttt{r}, \texttt{sanitise(r) $\in$ accepted} --- the model's guard discharged as
a totality obligation on the sanitiser rather than as a runtime check on the
commit.  With that obligation met, \texttt{apply\_patch}'s boundary re-check
becomes a never-tripped invariant, the arbiter's forget branch always eventually
fires, and no rejected value can be honoured.

The obligation is exactly the one the parent model discharged \emph{structurally}
for the colour carrier --- a hex string cannot encode a non-finite value, so its
parse is total onto \texttt{[0,1]} --- and \emph{by validation} for the slider,
whose \texttt{validate} carries a \texttt{math.isfinite} guard because raw floats
cross the wire.  The numeric field is the slider's case: raw numbers cross the
wire, so finiteness must be enforced, and the natural place is the same
pre-commit projection that already enforces range.

% ============================================================
\section{Scope and Relation to the Parent Model}
\label{sec:scope}
% ============================================================

This specification isolates the \emph{convergence} concern and takes the parent
model's honour/defer/clobber concerns as given.  It keeps the edit cycle
(\texttt{Focus}, \texttt{Keystroke}, \texttt{Commit}, the two honour frames) so
that the convergence question is asked over a faithful commit-echo loop, but it
does not restate the durability (\texttt{lost}) or no-clobber
(\texttt{clobbered}) invariants; those are proved in
\texttt{input\_text\_reconciliation.tex} and are unaffected by the commit guard,
which only narrows the set of committable values.

Two modelling choices bound the claim.  First, \texttt{accepted} is static: the
Hub's accepted domain does not change under the field's own operations, which
matches a field whose \texttt{min}/\texttt{max}/finiteness rule is fixed for the
life of the element.  A field whose bounds an agent narrows \emph{while a commit
is in flight} could reject a previously-acceptable commit; that is a distinct
interleaving, outside this model, and would be caught by the same invariant were
\texttt{accepted} made mutable.  Second, a single \texttt{committed} component
admits only one outstanding commit, as in the parent model; the honour-forever
fault needs only one, and modelling a queue would add interleaving the fault does
not require.  Both are consistent with the parent model's single-slot,
value-equality reconciliation.

\end{document}
